\section{Objectives}

\subsection{General Objective}
The primary objective of this thesis is to evaluate the suppression of the Saffman--Taylor instability in porous media during Enhanced Oil Recovery (EOR) processes using magnetic fluids. This is achieved through the development of an integrated framework comprising a rigorous analytical mathematical formulation and a Computational Fluid Dynamics (CFD) simulation.

\subsection{Specific Objectives}
To accomplish these goals, the following specific objectives are established:

\begin{itemize}
    \item \textbf{Analytical Formulation:} Formulate a Linear Stability Analysis (LSA) for the multiphase displacement involving a magnetic fluid under an external magnetic field, in order to derive the theoretical stability criteria and the perturbation dispersion relations.

    \item \textbf{Computational Modeling:} Develop and verify a multiphase CFD model based on the Cahn--Hilliard phase-field formulation coupled to a Lattice Boltzmann hydrodynamic solver, capable of simulating the nonlinear displacement of heavy oil by a magnetic fluid within a representative porous medium geometry.

    \item \textbf{Theoretical Validation:} Compare and validate the numerical predictions of the CFD model against the analytical thresholds established by the LSA, ensuring the physical accuracy of the computational framework prior to complex extrapolations.

    \item \textbf{Interfacial Dynamics Analysis:} Investigate the local fluid--fluid interaction, quantifying how varying magnetic field topologies physically alter the interfacial dynamics, either mitigating or potentially exacerbating finger growth.

    \item \textbf{Regime Mapping and EOR Evaluation:} Conduct a parametric study to map the flow regimes, comparing baseline conditions (conventional oil--water displacement) with magnetic fluid injection.
\end{itemize}